% ===================================================================
%  ТАБЛИЦА № 8 — Жатва. — Охота, Сбор винограда, Рыбная ловля.
%  Traduction 2026 d'apres texto/io, controlee sur texto/fr, texto/en
%  et texto/es. Consigne : voir texto/ru/10-tablica-01.tex.
%
%  L'appel de note est dans le TITRE de la scene — « Жатва (*) » — et
%  non dans un alinea : c'est ainsi que l'ido le compose.
% ===================================================================

%%K t08-noto-1 noto
\VUnotes{143.2mm}{%
\textit{(*) Потому что это слово неточно: жатва льна, сбор винограда,
сбор яблок? Быть может, лучше было бы взять «} \VUgras{mesono}
\textit{», предложенное в} Mondo \textit{XIV, стр. 242. Но пока этот
новый корень не принят Академией и ждёт обсуждения по обычным правилам
идо-движения.}%
}

%%K t08-apar-1 apar
\VUsaut{5.0mm}
\VUtitre{15.0pt}{60}{ТАБЛИЦА № 8}
\VUsaut{3.0mm}
\VUcentre{13.2pt}{90}{\textit{Первая сцена.}}
\VUsaut{3.0mm}
\VUpk{10.2pt}{40}{Жатва (*).}
\VUsaut{4.0mm}
\VUpk{10.2pt}{40}{Лицо полей.}
\VUsaut{3.0mm}

%%K t08-c1-01-1 p
1. --- С приходом \VUgras{лета} поля снова переменили лицо. Семя, вверенное
борозде, взошло; из земли поднялось маленькое зелёное растение и дало
свой колос. Зерно налилось, потом созрело, и теперь остаётся только
убрать жатву.

%%K t08-c1-02-1 p
2. --- \VUgras{Полевые цветы} распустились. \VUgras{Васильки}
\textsuperscript{(1)}, \VUgras{маки} \textsuperscript{(2)} и \VUgras{наперстянки}
\textsuperscript{(3)} вносят в ровный тон хлебов весёлый спор своих свежих
\VUgras{венчиков}.

%%K t08-c1-03-1 p
3. --- Плодовые деревья оделись листвой; плоды поспели. Маленькая
\VUgras{вишня} \textsuperscript{(4)} гнётся под прекрасными \VUgras{вишнями}
\textsuperscript{(5)}, которые дети рвут и едят с наслаждением.

%%K t08-c1-04-1 p
4. --- Стадо может теперь целый день быть на воле.

%%K t08-c1-04-2 p
\VUgras{Ягнята} \textsuperscript{(6)} выросли и стали славными \VUgras{овцами}
\textsuperscript{(7)} и славными \VUgras{баранами} \textsuperscript{(8)} с густым
руном. Они мирно щиплют \VUgras{траву} на \VUgras{пастбище}
\textsuperscript{(9)}, усеянном \VUgras{маргаритками}.

%%K t08-c1-04-3 p
Скоро с этих животных состригут \VUgras{шерсть}, из которой делают сукно
для нашей одежды.

%%K t08-tit-2 sub
\VUsaut{5.0mm}
\VUpk{10.2pt}{40}{Пастух.}
\VUsaut{3.4mm}

%%K t08-c1-05-1 p
5. --- \VUgras{Пастух} \textsuperscript{(10)}, кажется, не очень-то на них
смотрит. Его занимает только гроза, проходящая на горизонте, а
\VUgras{овчар} \textsuperscript{(13)}, поднося ко рту свою \VUgras{флягу}
\textsuperscript{(14)}, кажется ещё беззаботнее.

%%K t08-c1-06-1 p
6. --- Жара давит; вот и \VUgras{собака} \textsuperscript{(15)} идёт освежиться;
она лакает свежую воду \VUgras{ручья} \textsuperscript{(16)} и пугает бедную
маленькую \VUgras{лягушку} \textsuperscript{(17)}, притаившуюся в траве на
берегу. Та прыгает в прозрачную воду, где цветут \VUgras{кувшинки}
\textsuperscript{(18)}.

%%K t08-c1-07-1 p
7. --- \VUgras{Трясогузка} \textsuperscript{(19)} купается возле маленького
\VUgras{родника} \textsuperscript{(20)}, который бьёт между двух камней и прячется
под \VUgras{колючим кустарником} \textsuperscript{(21)} и \VUgras{кустами
боярышника} \textsuperscript{(22)}.

%%K t08-tit-3 sub
\VUsaut{5.0mm}
\VUpk{10.2pt}{40}{Жатва.}
\VUsaut{3.4mm}

%%K t08-c1-08-1 p
8. --- Для деревенских детей уборка хлеба --- время очень весёлое. Они
весело \VUgras{кувыркаются} \textsuperscript{(24)} на \VUgras{копнах соломы}
\textsuperscript{(23)}, помогают \VUgras{жнецам} \textsuperscript{(26)}, и пока те
косят \VUgras{ниву} \textsuperscript{(27)} своими \VUgras{косами} \textsuperscript{(25)},
хорошо отбитыми на \VUgras{бруске} \textsuperscript{(30)}, они собирают хлеб и
вяжут его в \VUgras{снопы} \textsuperscript{(31)} или \VUgras{вязанки}
\textsuperscript{(32)}.

%%K t08-c1-09-1 p
9. --- Иногда самым большим дают срезать \VUgras{серпом} \textsuperscript{(12)}
\VUgras{золотые стебли} \textsuperscript{(28)}, несущие тяжёлые \VUgras{колосья}
\textsuperscript{(29)}, полные зерна; а маленькие, присматривая за
\VUgras{скотиной} \textsuperscript{(33)}, что пасётся на \VUgras{лугу}
\textsuperscript{(34)} в тени большой \VUgras{липы} \textsuperscript{(35)}, ловят
своим \VUgras{сачком} \textsuperscript{(36)} красивых \VUgras{бабочек}.

%%K t08-c1-09-2 p
Иные даже влезают на \VUgras{воз} \textsuperscript{(37)}, чтобы уложить снопы,
которые им подают на \VUgras{вилах} \textsuperscript{(38)}.

%%K t08-c1-10-1 p
10. --- По всей \VUgras{равнине} \textsuperscript{(39)}, где взлетают
\VUgras{жаворонки} \textsuperscript{(41)}, идёт большая суета. Все заняты жатвой.
Но пока бедняки по-прежнему берутся за старые орудия --- \VUgras{косы,
серпы} и \VUgras{цепы} \textsuperscript{(40)} --- богатые землевладельцы жнут
свою пшеницу или свою \VUgras{рожь} \textsuperscript{(43)} \VUgras{жатвенной
машиной} \textsuperscript{(42)}. Затем, не возя снопов в амбар, тут же на месте
складывают большие \VUgras{скирды} \textsuperscript{(44)}.

%%K t08-c1-11-1 p
11. --- Потом подвозят \VUgras{молотилку} \textsuperscript{(47)}, которую вертит
\VUgras{локомобиль} \textsuperscript{(45)} через \VUgras{приводной ремень}
\textsuperscript{(46)}, и так зерно отделяется от \VUgras{соломы}, не покидая
поля, где было посеяно.

%%K t08-tit-4 sub
\VUsaut{5.0mm}
\VUpk{10.2pt}{40}{Гроза.}
\VUsaut{3.4mm}

%%K t08-c1-12-1 p
12. --- Во время жатвы часто разражаются сильные \VUgras{грозы}
\textsuperscript{(53)}. Сперва издали слышится раскат \VUgras{грома}; поднимается
сильный \VUgras{западный ветер} \textsuperscript{(54)} и до стона гнёт самые
большие деревья \VUgras{леса} \textsuperscript{(49)}. Чёрные \VUgras{тучи}
\textsuperscript{(56)} берут небо приступом; их пересекают слепящие зигзаги
\VUgras{молнии} \textsuperscript{(55)}. Начинает падать \VUgras{дождь}, а порой и
\VUgras{град}, прибивая готовый к жатве хлеб.

%%K t08-c1-13-1 p
13. --- Часто молния зажигает какое-нибудь строение. Так, в прошлом году
крыша \VUgras{ветряной мельницы} \textsuperscript{(50)} сгорела дотла, а два её
\VUgras{крыла} \textsuperscript{(51)} были разбиты.

%%K t08-c1-14-1 p
14. --- Через несколько часов возвращается покой. На горизонте является
блистательная \VUgras{радуга} \textsuperscript{(52)}, и природа снова берётся за
жизнь, прерванную на время. Деревенские, укрывшиеся в \VUgras{шалашах}
\textsuperscript{(48)}, возвращаются к работе и храбро принимаются поправлять
только что понесённый урон.

%%K t08-tit-5 sub
\VUsaut{6.0mm}
\VUcentre{13.2pt}{90}{\textit{Вторая сцена.}}
\VUsaut{3.6mm}

%%K t08-tit-6 sub
\VUpk{10.2pt}{40}{Охота. --- Сбор винограда.}
\VUsaut{1.4mm}
\VUpk{10.2pt}{40}{Рыбная ловля.}
\VUsaut{4.0mm}

%%K t08-c2-01-1 p
1. --- К концу \VUgras{августа} или к середине \VUgras{сентября}, смотря по
краю, открывается охота. Едва первые лучи солнца выманят \VUgras{ящериц}
\textsuperscript{(2)} из нор, как \VUgras{охотник} \textsuperscript{(5)} уже выходит
со своими \VUgras{собаками} \textsuperscript{(4)}.

%%K t08-c2-02-1 p
2. --- Он надел свой прочный \VUgras{охотничий костюм} \textsuperscript{(9)} из
толстого сукна и кожаные \VUgras{гетры}, которые позволят ему идти по
мокрой траве, где прячутся \VUgras{жабы} \textsuperscript{(1)}; он взял
\VUgras{ружьё} \textsuperscript{(6)} и \VUgras{ягдташ} \textsuperscript{(10)} --- и в
путь.

%%K t08-c2-03-1 p
3. --- Азарт охоты приводит его в самую середину виноградника, где сбор в
полном разгаре. Дети виноградаря, которые обыкновенно стерегут коз у
дороги, сегодня пускают их пастись где вздумается и, привязав к колу
старого \VUgras{козла} \textsuperscript{(3)}, который наверняка воспользовался бы
свободой, чтобы удрать, тоже берут свою долю веселья. Один достаёт гроздь
из \VUgras{корзины для винограда} \textsuperscript{(12)} и забавляется, выжимая
\VUgras{сок} \textsuperscript{(14)} в \VUgras{чашку} \textsuperscript{(13)}, чтобы
получилось сусло (неперебродившее вино). Другой венчает себя
\VUgras{венком из листьев} \textsuperscript{(15)}, который только что сплёл. Возле
них дерзкие \VUgras{воробьи} \textsuperscript{(16)} подлетают к самому прессу,
чтобы стащить виноград.

%%K t08-c2-04-1 p
4. --- Пока дети забавляются, мужчины работают. \VUgras{Сборщики}
\textsuperscript{(18)} носят виноград в подвал в \VUgras{заплечных корзинах}
\textsuperscript{(17)}; \VUgras{бондарь} \textsuperscript{(19)} чинит \VUgras{бочки}
\textsuperscript{(20)}, проверяет, целы ли \VUgras{клёпки} \textsuperscript{(22)} и
\VUgras{донья} \textsuperscript{(21)}, и заменяет лопнувшие \VUgras{обручи}
\textsuperscript{(23)}. Это он же сделал большие \VUgras{чаны} \textsuperscript{(25)},
куда ссыпают \VUgras{виноград} \textsuperscript{(26)}, чтобы он бродил.

%%K t08-c2-05-1 p
5. --- Когда брожение кончится, вино сливают, а остаток кладут под
\VUgras{пресс} \textsuperscript{(28)}; понемногу завинчивая большой \VUgras{винт}
\textsuperscript{(29)}, выжимают сок, который стекает в \VUgras{чан}
\textsuperscript{(32)} и даёт ещё \VUgras{вина} \textsuperscript{(31)}.

%%K t08-tit-8 sub
\VUsaut{6.0mm}
\VUpk{10.2pt}{40}{Пиво и сидр.}
\VUsaut{3.4mm}

%%K t08-c2-06-1 p
6. --- По стене \VUgras{погреба} \textsuperscript{(27)} вьётся ветка
\VUgras{хмеля} \textsuperscript{(30)}. Здесь это лишь украшение, но в иных
странах, где лоза очень редка, хмель разводят, чтобы делать \VUgras{пиво}.

%%K t08-c2-07-1 p
7. --- Маленькая \VUgras{яблоня} \textsuperscript{(33)} гнётся под яблоками,
которые дадут отличный \VUgras{сидр}, когда созреют. В своей \VUgras{клетке}
\textsuperscript{(34)} \VUgras{канарейка} \textsuperscript{(35)} и \VUgras{щегол}
\textsuperscript{(36)} завистливо глядят на воробьёв, резвящихся на воле.

%%K t08-c2-08-1 p
8. --- Насколько хватает глаз, по склону холмов стоят ряды
\VUgras{тычин} \textsuperscript{(40)}, поддерживающих великолепные \VUgras{лозы}
\textsuperscript{(39)}, тяжёлые от \VUgras{гроздьев}.

%%K t08-c2-08-2 p
Круглый год \VUgras{виноградарь} \textsuperscript{(42)} трудился, чтобы ходить за
своим \VUgras{виноградником} \textsuperscript{(38)}, и вот он вознаграждён за
труды хорошим урожаем. \VUgras{Сборщицы} \textsuperscript{(41)} наполняют чаны,
поставленные на воз, который тянут \VUgras{мулы} \textsuperscript{(43)}, и все
веселы.

%%K t08-tit-9 sub
\VUsaut{5.0mm}
\VUpk{10.2pt}{40}{Рыбная ловля.}
\VUsaut{3.4mm}

%%K t08-c2-09-1 p
9. --- То же и у \VUgras{удильщиков} \textsuperscript{(44)}, которые пришли на
рассвете расположиться у воды со своими \VUgras{удочками} \textsuperscript{(45)}.

%%K t08-c2-09-2 p
\VUgras{Рыба} \textsuperscript{(47)} берёт крючок, едва они опустят в воду свою
\VUgras{леску} \textsuperscript{(46)}, и они едва успевают сменить \VUgras{насадку},
как новая жертва уже бьётся на конце лески.

%%K t08-c2-09-3 p
Всё же \VUgras{рыбак с сетью} \textsuperscript{(48)}, который бросает свою сеть с
\VUgras{лодки} \textsuperscript{(49)} на середину \VUgras{реки} \textsuperscript{(50)},
ловит куда больше.

%%K t08-c2-10-1 p
10. --- Осень --- пора хороших прогулок: пешком, \VUgras{верхом}
\textsuperscript{(52)}, на \VUgras{велосипеде} \textsuperscript{(53)}, на
\VUgras{автомобиле} \textsuperscript{(54)}. \VUgras{Дороги} \textsuperscript{(51)}
чисты, и люди пользуются последними погожими днями, потому что
\VUgras{ласточки} \textsuperscript{(56)} уже собираются на крышах, чтобы улететь
в тёплые края. Листья старого \VUgras{орешника} \textsuperscript{(57)} желтеют,
\VUgras{орехи} падают, а высокие \VUgras{деревья}, стоящие \VUgras{рощей}
\textsuperscript{(58)} на \VUgras{возвышенности} \textsuperscript{(59)}, скоро
оголятся.

%%K t08-c2-11-1 p
11. --- \VUgras{Солнце} \textsuperscript{(60)} встаёт позже, дни делаются короче,
\VUgras{поутру} начинает чувствоваться довольно резкий холодок, и уже
стаи \VUgras{вальдшнепов} \textsuperscript{(61)} покидают наш негостеприимный
край.
\VUsaut{6.0mm}
\VUfilet{22mm}
